\chapter*{ВВЕДЕНИЕ}
\addcontentsline{toc}{chapter}{ВВЕДЕНИЕ}
\sloppy

\textbf{Актуальность темы исследования.}
Большие языковые модели (LLM) всё активнее встраиваются в системы поддержки принятия решений: от медицинской диагностики и юридического анализа до финансового консультирования и автоматизированного обучения~\cite{bommasani2021opportunities}. По мере того как область применения подобных систем расширяется, принципиальное значение приобретает вопрос об их когнитивной надёжности: способны ли модели рассуждать нормативно корректно, или же они воспроизводят систематические ошибки, характерные для человеческого мышления? Когнитивные предубеждения~--- устойчивые отклонения от нормативных стандартов вероятностного и логического рассуждения~--- давно являются предметом исследований в психологии и поведенческой экономике~\cite{kahneman2011thinking}. С появлением масштабных языковых моделей, обученных на огромных корпусах текстов, порождённых людьми, закономерно возникает гипотеза о том, что подобные предубеждения могут быть имплицитно усвоены в ходе обучения и проявляться при решении задач, требующих строгого вероятностного или логического мышления.

Несмотря на растущий массив литературы, посвящённой рациональности LLM~\cite{binz2023using, suri2024large, macmillan2024irrationality}, большинство существующих работ сосредоточено на отдельных задачах или узком наборе моделей, не рассматривая предубеждения как феномен, существующий на спектре~--- от вероятностного суждения к классическому логическому рассуждению и далее к индуктивному открытию правил. Настоящая работа восполняет этот пробел, предлагая систематическое сравнительное исследование трёх парадигм когнитивной психологии на девяти современных языковых моделях.

\textbf{Степень разработанности проблемы.}
Первые систематические исследования когнитивных предубеждений в LLM появились в 2023--2024 годах. Binz и Schulz~\cite{binz2023using} показали, что GPT-3 демонстрирует эффекты репрезентативности и предвзятости подтверждения на классических задачах из программы эвристик и предубеждений Канемана и Тверски. Suri и соавторы~\cite{suri2024large} расширили эту линию, исследовав эффект фрейминга и ошибку конъюнкции в GPT-4. Macmillan-Scott и Musolesi~\cite{macmillan2024irrationality} провели масштабный бенчмарк нескольких моделей на задаче Уэйсона. Тем не менее ни одна из указанных работ не рассматривает все три парадигмы~--- конъюнктивную ошибку, задачу выбора карточек и задачу открытия правила 2-4-6~--- в едином экспериментальном дизайне и не вводит непрерывных метрик, позволяющих измерить силу предубеждения за пределами бинарной классификации ответов.

\textbf{Цель и задачи исследования.}
Цель настоящей работы~--- систематически измерить и сравнить когнитивные предубеждения в современных LLM на трёх классических психологических парадигмах, охватывающих различные типы рассуждения: вероятностное суждение, условная логика и индуктивное открытие правил.

Для достижения данной цели поставлены следующие задачи:
\begin{itemize}
  \item провести анализ существующих подходов к изучению когнитивных предубеждений в LLM и выявить методологические ограничения предшествующих работ;
  \item разработать экспериментальный дизайн для трёх парадигм~--- задачи Линды (ошибка конъюнкции), задачи выбора карточек Уэйсона и задачи открытия правила 2-4-6~--- с контролем за контаминацией обучающих данных;
  \item сформировать оригинальные датасеты для каждого эксперимента, обеспечивающие структурное единообразие и отсутствие пересечения с обучающими корпусами;
  \item провести эксперименты на девяти современных языковых моделях, охватывающих широкий спектр архитектур, размеров и провайдеров;
  \item ввести и апробировать непрерывную метрику силы предубеждения (Severity of Bias) как альтернативу бинарной классификации ошибок;
  \item сравнить поведение моделей в трёх парадигмах и проверить гипотезу о масштабировании как факторе снижения предубеждений;
  \item на основании полученных результатов сформулировать рекомендации по выбору и оценке LLM для приложений, требующих нормативно корректного рассуждения.
\end{itemize}

\textbf{Объект и предмет исследования.}
Объектом исследования являются большие языковые модели как системы рассуждения. Предметом исследования служат когнитивные предубеждения в ответах LLM на задачи из трёх психологических парадигм: ошибка конъюнкции, предвзятость подтверждения в условной логике и предвзятость подтверждения в индуктивном открытии правил.

\textbf{Научная новизна.}
Научная новизна работы определяется следующими результатами:
\begin{itemize}
  \item впервые проведено сравнительное исследование трёх классических парадигм когнитивной психологии~--- ошибки конъюнкции, задачи Уэйсона и задачи 2-4-6~--- в едином экспериментальном дизайне на одном и том же наборе из девяти моделей, что позволяет рассматривать предубеждение как единый феномен на спектре типов рассуждения;
  \item введена оригинальная непрерывная метрика Severity of Bias, основанная на самооценке уверенности модели и позволяющая обнаруживать токен-предубеждение в случаях, когда бинарная метрика демонстрирует потолок точности;
  \item разработан оригинальный датасет из 110 биографических виньеток для задачи Линды с систематической факториальной подстановкой демографических переменных, что обеспечивает контроль за конфаундерами нарративного стиля;
  \item введено условие «фантастических социальных контрактов» в задаче Уэйсона, позволяющее диссоциировать эффект знакомости содержания от эффекта деонтической структуры, что, насколько известно автору, ранее не реализовывалось в исследованиях LLM;
  \item обнаружен и описан внутрисемейный парадокс масштаба в семействах моделей Gemma и Qwen: меньшая модель демонстрирует более низкий уровень предубеждения, чем большая модель того же семейства.
\end{itemize}

\textbf{Практическая значимость.}
Результаты исследования имеют непосредственное прикладное значение в нескольких направлениях. Во-первых, разработанный фреймворк тестирования когнитивных предубеждений может применяться при выборе языковой модели для систем поддержки принятия решений, где критична нормативная корректность рассуждения~--- в частности, в медицинской диагностике, юридическом анализе и финансовом консультировании. Во-вторых, введённая метрика Severity of Bias предоставляет более тонкий инструмент оценки, чем бинарная точность, что позволяет проводить сравнение моделей в условиях, когда все они достигают высокой точности по бинарной метрике. В-третьих, выявленные различия в поведении моделей на различных типах задач позволяют формулировать рекомендации по prompt engineering для снижения предубеждений. Полученные датасеты и экспериментальный код публикуются в открытом доступе в целях воспроизводимости исследований.

\textbf{Методологическая основа.}
Исследование опирается на классические парадигмы когнитивной психологии~\cite{tversky1983extensional, wason1966reasoning, wason1960failure}, адаптированные для тестирования языковых моделей. В качестве статистических инструментов используются точный тест МакНемара для согласованных пар~\cite{mcnemar1947note}, парный t-тест, критерий Уилкоксона и процедура Бенджамини-Хохберга для контроля частоты ложных открытий~\cite{benjamini1995controlling}. Для оценки качества гипотез в задаче 2-4-6 применяется метрика расширенного пересечения по объединению (IOU) на основе метода Монте-Карло, впервые предложенная Banatt и соавторами~\cite{banatt2024wilt}.

\textbf{Апробация результатов.}
Основные результаты исследования были представлены на кандидатском экзамене в AI Talent Hub ИТМО (январь 2026~г.). Работа выполнена в рамках научно-исследовательской практики магистерской программы по направлению искусственный интеллект под руководством Ильи Макарова.

\textbf{Структура работы.}
Работа состоит из введения, трёх глав, заключения, списка использованных источников и приложений.

В \textit{первой главе} проводится анализ предметной области: рассматриваются теоретические основы когнитивных предубеждений, анализируется проблема их воспроизведения в LLM в соответствии с текущим состоянием, выполняется обзор смежных работ и формулируются исследовательские гипотезы.

Во \textit{второй главе} описывается дизайн экспериментов: обосновывается выбор парадигм, представлены оригинальные датасеты, промпт-условия, протестированные модели и система метрик; приводятся результаты экспериментов и их анализ.

В \textit{третьей главе} предлагается архитектура практического инструмента для оценки когнитивных предубеждений в LLM на основе полученных экспериментальных результатов.

В \textit{заключении} подводятся итоги работы, формулируются основные выводы и намечаются направления дальнейших исследований.